%-------------------------------------------------------------------------------
\section{Cohesive sediment transport}
%-------------------------------------------------------------------------------
Cohesive properties appear for fine particles (silts and clay), with diameter less than a limiting value of about 60 $\mu$m, depending on the physico-chemical
properties of the fluid and salinity. The separation value at $60\mu$m to discriminate non-cohesive from cohesive sediment is conventional. This value is different depending on the country (e.g. $63\mu$m in The Netherlands, $75\mu$m in USA as pointed by Winterwerp and Van Kesteren~\cite{Winterwerp}). Moreover, aggregation of flocs can lead to the formation of macro-flocs larger than $100\mu$m.\\

Fine cohesive sediments are mainly transported in suspension and transport processes strongly depend on the state of flocculation of the suspension and consolidation of the bed. The erosion rate mainly depends on the degree of consolidation of the sediment bed, while the settling velocity depends on the state of flocculation and aggregates properties.\\

In \gaia{}, cohesive sediments are accounted by solving the 2D advection-diffusion equation:
\begin{equation*}
\frac{\partial hC}{\partial t} + \frac{\partial hUC}{\partial x} + \frac{\partial hVC}{\partial y} =
\frac{\partial}{\partial x}\left(h\epsilon_s\frac{\partial C}{\partial x}\right) +
\frac{\partial}{\partial y}\left(h\epsilon_s\frac{\partial C}{\partial y}\right) + (E-D)
\end{equation*}
$C=C(x,y,t)$ is the depth-averaged concentration \textcolor{black}{\bf expressed in g/l}, $(U,V)$ are the depth-averaged components of the velocity in the $x$ and $y$ directions, respectively, $\epsilon_s$ is the turbulent diffusivity of the sediment. Variables $E$ and $D$ are respectively the erosion and deposition fluxes.
